\section{结论}

本文研究了递归历史在镜头切换处的冲突作用：切换前历史为建立连续镜头之间的关系提供参照，但将同一状态继续传播到新镜头可能引起随时间变化的镜头内漂移，而单独重新初始化又会丢失跨镜头参照。\method{} 只在临时对齐路径中使用历史，独立重新初始化的状态则负责后续递归。基于预测几何的人物关联连接跨镜头身份，共享变换将相机、人体和场景预测映射到统一世界坐标系。

受控实验支持将对齐与状态传播视为两个独立决策。在三个基准上，相比连续状态递归，\method{} 均改善了 WA-MPJPE、相机轨迹精度和 IDF1，并在大视角变化下获得最明显的收益。这些改进无需未来帧或逐视频优化；在包含一次镜头切换的标准化重建测试中，相比 \humanthree{} 仅增加 3.8\% 的时间。更进一步，我们的结果表明，递归历史可以作为镜头间对齐的临时坐标参照，而不必作为新镜头的递归状态继续传播。复杂剪辑、人物反复进出画面以及快速视觉变化仍是开放问题。
